% BEGIN CR28_CONSTITUTIVE_INVERSE
\section{Récupération fonctionnelle, composition et moment de Catalan}
\label{cr28:sec:restitucion-constitutiva}

La réponse des incidences \(90/120\) conserve plus d'information qu'une valeur scalaire isolée. Dans la même notation que \eqref{eq:three-four-completion}, écrivons
\[
 f(s)=\frac{1+s+s^2}{1+s+s^2+s^3},\qquad
 \mathcal D=s\,\frac{d}{ds},\qquad
 \mathcal V=f(T^{30}).
\]
Les marques \(P_\pm\) maintiennent l'orientation des canaux. Le transport angulaire est l'antécédent de la réponse ; sa reconstruction fonctionnelle retrouve cette représentation sans remplacer l'archive de mémoire dont elle provient.

\begin{theorem}[Inversion de la réponse constitutive]
\label{cr28:thm:inversion-constitutiva} La fonction \(f\) est une bijection décroissante de \((0,1)\) sur \((3/4,1)\). Son inverse \(\psi\) associe à \(r\) l'unique racine \(s\in(0,1)\) de
\begin{equation}
 r s^3+(r-1)(1+s+s^2)=0.
 \label{cr28:eq:cubic}
\end{equation}
Les deux réponses étiquetées \(r_\pm\) déterminent
\begin{equation}
 s_\pm=\psi(r_\pm),\quad q_\pm=s_\pm^{1/30},\quad
 x=-\frac12\log(q_+q_-),\quad
 y=\frac12\log\frac{q_-}{q_+}.
 \label{cr28:eq:angular-recovery}
\end{equation}
Pour l'opérateur complet \(S=\mathcal V^2\), positif et de spectre dans \((9/16,1)\), le calcul fonctionnel fournit
\begin{equation}
 T=\bigl[\psi(S^{1/2})\bigr]^{1/30},\qquad
 \det T=\det\bigl[\psi(S^{1/2})\bigr]^{1/30}.
 \label{cr28:eq:operator-recovery}
\end{equation}
\end{theorem}
\begin{proof}
La dérivée négative de \eqref{eq:vacuum-monotonic34} et les limites \(f(0^+)=1,\ f(1^-)=3/4\) prouvent la bijection. Multiplier \(f(s)=r\) par son dénominateur positif donne \eqref{cr28:eq:cubic}. La racine positive d'ordre \(30\) retrouve \(q_\pm\) ; la somme et la différence de \(\log q_\pm=-x\mp y\) donnent les deux coordonnées angulaires. Par positivité, \(S^{1/2}=\mathcal V\). Appliquer \(\psi\) à chaque projecteur spectral retrouve \(T^{30}\) et la racine positive donne \(T\). Le calcul de son déterminant conclut.
\end{proof}

Cette récupération de l'opérateur complet ne s'étend pas à un déterminant isolé. En effet,
\[
 T_1=\operatorname{diag}(1/5,1/3),\qquad
 T_2=\operatorname{diag}(1/6,2/5)
\]
ont pour déterminant \(1/15\), mais possèdent des spectres distincts dans la même chambre contractive. L'injectivité démontrée implique \(f(T_1^{30})^2\ne f(T_2^{30})^2\). Les deux canaux et leurs marques préservent la récupération ; leur réduction à un produit perd de l'information.
% END CR28_CONSTITUTIVE_INVERSE

% BEGIN CR28_TRANSPORT_COMPOSITION
% Recuperación expositiva de INTERRELACIONES_ESTRUCTURALES_VACIO.md, §§6--8.
% Insertar después de la inversión cúbica y antes de las coordenadas logarítmicas.
\subsection{Composition transportée des réponses constitutives}
\label{cr28:comp-sec}

Le transport angulaire et les incidences $90$ et $120$ proviennent de la
même structure APP--TRIT--TPK : les feuillets conservent les évaluations et
leurs retenues ; le régime tritique conserve l’orientation ; le calendrier
TPK transporte la phase et la mémoire jusqu’aux canaux angulaires. La réponse
constitutive applique à ces canaux la fonction $f$ de
\eqref{eq:vacuum-monotonic34}. Son inversion permet d’exprimer, dans les coordonnées
de réponse, la composition des transports du même repère. On conserve
ainsi l’ordre constructif : transport, composition et évaluation spectrale.
L’opération obtenue agit sur cette représentation de l’état enrichi ;
son inverse retrouve les canaux, tandis que l’histoire complète de mémoire
demeure dans la structure de provenance.

Soient $s=q^{30}$ et $\psi=f^{-1}$. L’exposant $30$ est le plus grand commun
diviseur des incidences $90$ et $120$. L’inversion cubique détermine
$\psi:(3/4,1)\longrightarrow(0,1)$. En adjoignant la réponse de frontière
$f(1)=3/4$, nous étendons la bijection à
\[
 \psi:D\longrightarrow(0,1],\qquad D=[3/4,1).
\]
Sur cette frontière, on utilise la fonction rationnelle $f$ ; le quotient
$(1-q^{90})/(1-q^{120})$ y a une singularité éliminable.

\begin{proposition}[Monoïde des réponses et coordonnée additive]
\label{cr28:comp-monoid}
L’opération
\begin{equation}
 r\odot t=f\bigl(\psi(r)\psi(t)\bigr),\qquad r,t\in D,
 \label{cr28:comp-law}
\end{equation}
fait de $D$ un monoïde commutatif et simplifiable d’élément neutre $3/4$.
L’application
\begin{equation}
 \ell(r)=-\frac1{30}\log\psi(r)
 \label{cr28:comp-character}
\end{equation}
est un isomorphisme de ce monoïde avec $(\mathbb R_{\geq0},+)$.
L’intervalle $(3/4,1)$ est stable par $\odot$ ; le neutre appartient à
l’extension de frontière. Le seul élément inversible dans $D$
est le neutre.
\end{proposition}
\begin{proof}
Le produit de deux éléments de $(0,1]$ appartient au même intervalle et
\[
 \psi(r\odot t)=\psi(r)\psi(t).
\]
Pour trois réponses, les deux ordres d’association ont pour image
$\psi(r)\psi(t)\psi(u)$ ; l’injectivité de $\psi$ prouve
l’associativité. La commutativité provient du produit scalaire. Si
$r\odot t=r\odot u$, la positivité de $\psi(r)$ permet de simplifier et
d’obtenir $\psi(t)=\psi(u)$, donc $t=u$. Le paramètre $1$ correspond à
$3/4$ et fournit l’identité. Le logarithme transforme le produit en
somme, et son image dans $(0,1]$ prouve que $\ell$ est une bijection sur
$\mathbb R_{\geq0}$. En particulier, pour la réponse d’un canal $q$,
$\ell(r)=-\log q$. Si les deux paramètres sont strictement inférieurs à
$1$, leur produit l’est aussi. Enfin, deux nombres de $(0,1]$
ont pour produit $1$ uniquement s’ils sont tous deux $1$.
\end{proof}

\begin{proposition}[Simplification admissible et prolongement inverse]
\label{cr28:comp-inverse}
Pour $r,u\in D$, l’équation $r\odot t=u$ a une solution $t\in D$
exactement lorsque $u\geq r$. Cette solution est unique et donnée par
\begin{equation}
 t=f\!\left(\frac{\psi(u)}{\psi(r)}\right).
 \label{cr28:comp-cancellation}
\end{equation}
La même fonction $f$ est une bijection de $(0,\infty)$ sur $(0,1)$.
Avec cette extension de $\psi$, la loi~\eqref{cr28:comp-law} transporte le
groupe multiplicatif des paramètres positifs sur l’intervalle $(0,1)$.
L’inversion du groupe satisfait
\begin{equation}
 r^{\ominus}=\psi(r)r,\qquad
 r\odot r^{\ominus}=\frac34.
 \label{cr28:comp-group-inverse}
\end{equation}
\end{proposition}
\begin{proof}
L’équation équivaut à
$\psi(t)=\psi(u)/\psi(r)$. Ce quotient appartient à $(0,1]$
exactement lorsque $\psi(u)\leq\psi(r)$, ce qui, par monotonie décroissante,
équivaut à $u\geq r$. La bijection prouve l’existence et l’unicité.
La dérivée de \eqref{eq:vacuum-monotonic34} est négative pour tout $s>0$ ;
les limites de $f$ en $0$ et à l’infini sont $1$ et $0$, respectivement.
Cela établit la bijection étendue. En multipliant le numérateur et le
dénominateur par $s^3$, on obtient
\[
 f(s^{-1})=\frac{s^3+s^2+s}{s^3+s^2+s+1}=s f(s).
\]
L’inverse multiplicatif de $s=\psi(r)$ a donc pour réponse
$\psi(r)r$. Sa composition avec $r$ correspond au paramètre $1$.
\end{proof}

Le prolongement à $s>1$, ou de manière équivalente à $q>1$, est l’extension
algébrique du lecteur. Le secteur physique contractant du chapitre conserve
son domaine déclaré. La coordonnée $\ell$ se prolonge en un isomorphisme
avec $(\mathbb R,+)$ et change de signe sous l’inversion du groupe.

\subsubsection{Réalisation opératorielle dans un repère de feuillets commun}

\begin{proposition}[Composition avec projecteurs marqués]
\label{cr28:comp-operator}
Soient $\mathcal R=\mathcal R^*$, $\mathcal R^2=I$ et
$P_\pm=(I\pm\mathcal R)/2$. Pour $j=a,b$, soient
\[
 T_j=q_{j,+}P_++q_{j,-}P_-,\qquad
 \mathcal V_j=f(T_j^{30}),\qquad 0<q_{j,\pm}<1.
\]
La composition $T_{a\circ b}=T_aT_b$ satisfait
\begin{equation}
 \mathcal V_{a\circ b}
 =f\bigl(\psi(\mathcal V_a)\psi(\mathcal V_b)\bigr).
 \label{cr28:comp-operator-law}
\end{equation}
Ses deux réponses sont $r_{a,\pm}\odot r_{b,\pm}$. Si
$T_j=\exp[-(x_jI+y_j\mathcal R)]$ avec $x_j>|y_j|$, alors
\begin{equation}
 T_aT_b=
 \exp[-((x_a+x_b)I+(y_a+y_b)\mathcal R)],
 \label{cr28:comp-angular-addition}
\end{equation}
y $x_a+x_b>|y_a+y_b|$.
\end{proposition}
\begin{proof}
L’orthogonalité $P_+P_-=0$ donne
\[
 T_aT_b=q_{a,+}q_{b,+}P_++q_{a,-}q_{b,-}P_-.
\]
En particulier, $T_a$ et $T_b$ commutent. Le calcul fonctionnel dans ces
projecteurs fournit
$\psi(\mathcal V_j)=T_j^{30}$ et
$(T_aT_b)^{30}=T_a^{30}T_b^{30}$ ; appliquer $f$ prouve
\eqref{cr28:comp-operator-law}. Pour toute fonction $g$ définie sur les
deux valeurs propres, on a, plus explicitement,
\[
 P_\pm g(T_j)=g(q_{j,\pm})P_\pm.
\]
Comme $q_{j,\pm}=\exp[-(x_j\pm y_j)]$, leurs produits prouvent
\eqref{cr28:comp-angular-addition}. L’inégalité triangulaire conclut
$x_a+x_b>|y_a|+|y_b|\geq|y_a+y_b|$.
\end{proof}

Les coordonnées retrouvées peuvent s’écrire directement comme
\[
 x=\frac{\ell(r_+)+\ell(r_-)}2,\qquad
 y=\frac{\ell(r_+)-\ell(r_-)}2.
\]
La composition additionne ces paires. L’échange simultané de feuillets
permute les deux canaux de chaque transport, agit comme
$(x,y)\mapsto(x,-y)$ et est un automorphisme de la loi composante par
composante. L’inversion des deux transports dans l’extension algébrique
agit comme $(x,y)\mapsto(-x,-y)$. Les deux involutions commutent et
conservent des fonctions distinctes : l’une échange les marques d’orientation ;
l’autre inverse le transport.

Pour représenter l’échange par conjugaison, on utilise un
opérateur $L$ satisfaisant $L\mathcal R L^{-1}=-\mathcal R$. Dans la
représentation
\[
 \mathcal R=\begin{pmatrix}0&1\\1&0\end{pmatrix},\qquad
 L=\begin{pmatrix}1&0\\0&-1\end{pmatrix},
\]
la multiplication directe donne $LP_\pm L^{-1}=P_\mp$ et, par conséquent,
$LT_jL^{-1}=q_{j,-}P_++q_{j,+}P_-$. La conjugaison par
$\mathcal R$ conserve ses propres projecteurs. Une trace qui réunit les
canaux ne remplace pas non plus les projecteurs dans cette identité fonctionnelle :
l’information de leurs étiquettes intervient dans la composition.

L’hypothèse de repère commun est substantielle. Par exemple, les opérateurs
positifs définis
\[
 A_0=\begin{pmatrix}1/2&0\\0&1/3\end{pmatrix},\qquad
 B_0=\begin{pmatrix}1/2&1/6\\1/6&1/2\end{pmatrix}
\]
ont des valeurs propres positives, mais
\[
 A_0B_0=\begin{pmatrix}1/4&1/12\\1/18&1/6\end{pmatrix}
\]
n’est pas autoadjoint. La proposition spécifie la composition des
transports dans les projecteurs communs ; son domaine ne comprend pas un mélange
arbitraire de milieux matériels. La réalisation d’une concaténation de
chemins HMT conserve en outre ses données de mémoire : la fermeture algébrique du
cône de paramètres décrit la représentation et ne sélectionne pas à elle seule
de nouveaux chemins primaires.

\subsubsection{Produit de canaux et composition de transports}

La vitesse normalisée d’un état reste déterminée par
\[
 \widehat c=(r_+r_-)^{-1}.
\]
Ici, le produit ordinaire réunit les deux réponses de ce même état.
La loi $\odot$ évalue une autre opération : la composition de deux transports
dans chaque feuillet. Le calcul exact suivant distingue leurs arguments :
\begin{equation}
 f(1/2)=\frac{14}{15},\qquad
 \frac{14}{15}\odot\frac{14}{15}=f(1/4)=\frac{84}{85}
 \ne\frac{196}{225}=\left(\frac{14}{15}\right)^2.
 \label{cr28:comp-rational-example}
\end{equation}
Les fractions de cet exemple sont des paramètres algébriques de vérification.
L’expression de $\widehat c$ et les identités constitutives précédentes
restent inchangées. Le contenu ajouté est une loi explicite pour
transporter la composition et retrouver sa coordonnée angulaire additive.

% END CR28_TRANSPORT_COMPOSITION

% BEGIN CR28_CATALAN_RECOVERY
\subsection{Restitution fonctionnelle du moment orienté}
\label{cr28:sec:catalan-recovery}

La relation différentielle~\eqref{cr28:eq:catalan-problem} admet une inverse qui conserve le système complet des moments. Son domaine est la fonction constitutive produite par les secteurs du cycle, avec l'orientation de~\eqref{eq:vacuum-angular-channels}. Les incidences $90=3\cdot30$ et $120=4\cdot30$ fixent cette fonction avant son évaluation dans les canaux $s_\pm=q_\pm^{30}$. La condition à l'origine correspond à la disparition du moment lorsque son argument contractif s'annule. Cette condition détermine également les constantes d'intégration.

\begin{theorem}[Restitution intégrale et unicité]
\label{cr28:thm:catalan-recovery}
Soit $f(s)=(1-s^3)/(1-s^4)$, pour $0<s<1$, la réponse
de~\eqref{eq:three-four-completion}, et soit $\mathcal D=s\,d/ds$.
Il existe une unique fonction $F\in C^2((0,1))$ qui satisfait
\begin{equation}
 \mathcal D^2F(s)=\frac{s(1+s)}{1+s+s^2}f(s),
 \qquad \lim_{s\downarrow0}F(s)=0.
 \label{cr28:eq:catalan-problem}
\end{equation}
Elle est donnée par
\begin{align}
 F(s)&=\int_0^s\log\frac{s}{t}\,
          \frac{1+t}{1+t+t^2}f(t)\,dt
       =\int_0^s\frac{\log(s/t)}{1+t^2}\,dt
       =\mathcal C_2(s),
 \label{cr28:eq:catalan-integral}\\
 f(s)&=\frac{1+s+s^2}{s(1+s)}\mathcal D^2F(s).
 \label{cr28:eq:catalan-inverse}
\end{align}
L’extension continue à $s=1$ retrouve $F(1)=G_{\rm Cat}$.
\end{theorem}
\begin{proof}
La factorisation $1+t+t^2+t^3=(1+t)(1+t^2)$ donne
\[
 \frac{1+t}{1+t+t^2}f(t)=\frac1{1+t^2}.
\]
L’intégrande de~\eqref{cr28:eq:catalan-integral} est
intégrable en zéro, parce que
\[
 0\leq F(s)\leq\int_0^s\log(s/t)\,dt=s.
\]
Par conséquent, $F(s)\to0$. La différentiation sur l’intervalle ouvert,
d’abord sur un intervalle tronqué, puis par convergence dominée,
produit
\[
 \mathcal DF(s)=\int_0^s\frac{dt}{1+t^2},
 \qquad \mathcal D^2F(s)=\frac{s}{1+s^2}
     =\frac{s(1+s)}{1+s+s^2}f(s).
\]
En particulier, $F$ est de classe $C^2$ et satisfait l’équation.
Si $F_1,F_2$ sont deux solutions, la différence $H=F_1-F_2$
satisfait $\mathcal D^2H=0$. Au moyen de $z=\log s$, on obtient
$d^2H(e^z)/dz^2=0$, donc $H(s)=a\log s+b$.
L’existence de la limite nulle lorsque $s\downarrow0$ impose d’abord
$a=0$, puis $b=0$. La condition sur $F$ détermine donc
la solution et également la limite $\mathcal DF(s)\to0$.

Pour $s<1$, le développement géométrique de $(1+t^2)^{-1}$ et
l’intégrabilité absolue de ses termes pondérés permettent d’intégrer :
\[
 \int_0^s t^{2k}\log(s/t)\,dt
     =\frac{s^{2k+1}}{(2k+1)^2},\qquad
 F(s)=\sum_{k\geq0}\frac{(-1)^ks^{2k+1}}{(2k+1)^2}.
\]
La série représente le système intégral~\eqref{cr28:eq:catalan-integral}. La convergence absolue et uniforme de cette dernière série sur $[0,1]$ autorise le passage à l'extrémité $s=1$. Dans l'intégrale, la majoration intégrable $\log(1/t)$ l'autorise également, en prolongeant l'intégrande par zéro pour $t>s$. Enfin, isoler $f$ dans l'équation différentielle donne \eqref{cr28:eq:catalan-inverse}, avec des dénominateurs positifs sur $(0,1)$.
\end{proof}

La limite du moment et celle de sa dérivée seconde sont prises
dans leurs classes de convergence respectives. En particulier,
\[
 \lim_{s\uparrow1}\mathcal D^2\mathcal C_2(s)=\frac12
\]
est la limite radiale d’Abel de la série dérivée. L’évaluation
de $\mathcal C_2(1)$ utilise directement sa série absolument
convergente. Cette distinction conserve le même opérateur de profondeur
à l’intérieur et dans sa lecture de frontière.

\subsection{Compatibilité de la restitution avec les canaux et la vitesse}
\label{cr28:sec:channels-velocity}

Le théorème restitue une fonction à partir d'une autre fonction préalablement construite. L'inversion cubique~\eqref{cr28:eq:cubic} agit en revanche sur les évaluations $r_\pm$ : elle retrouve les arguments $s_\pm$ dans ce système fixé. La composition des deux opérations conserve l'orientation du quart de tour, les étiquettes $P_\pm$ et les incidences du transport. Ainsi, $G_{\rm Cat}=\mathcal C_2(1)$ est la lecture à l'extrémité d'un objet fonctionnel qui participe aussi aux réponses intérieures $r_\pm=f(s_\pm)$.

Pour $\mathcal V_m=\mathcal V\otimes I_{9^m}$ et
$\tau_m=\operatorname{Tr}/(2\cdot9^m)$, la lecture de vitesse
maintient exactement la normalisation de
\eqref{eq:vacuum-traces} :
\begin{equation}
 \exp[-2\tau_m(\log\mathcal V_m)]
   =\frac1{r_+r_-}
   =(\det\mathcal V)^{-1}=\widehat c.
 \label{cr28:eq:velocity-normalized}
\end{equation}
En effet, chaque valeur propre $r_\pm$ se répète $9^m$ fois, de sorte
que $2\tau_m(\log\mathcal V_m)=\log r_++\log r_-$.
L’identité utilise le déterminant de la réalisation à deux feuillets
et la trace normalisée de son amplification. Les sections
dimensionnelles déjà fixées donnent alors
\[
 c_{\rm int}=c_{\rm vac}=\widehat c\,u_C,
 \qquad \ell_0=c_{\rm vac}t_0,
\]
avec le même état, le même système de coordonnées et la même durée élémentaire que~\eqref{cr28:eq:same-velocity}. La restitution fonctionnelle conserve cette section de vitesse lorsqu'elle la compose avec l'action et la longueur ; la multiplicité spectrale est absorbée par la normalisation déclarée.

% END CR28_CATALAN_RECOVERY

% BEGIN CR28_CONSTITUTIVE_RADIAL_LINK
\subsection{Une même section constitutive dans l'action, la longueur et la gravité}
\label{cr28:sec:constitutive-radial} Les relations \eqref{eq:vacuum-sections} et la restitution précédente maintiennent une unique section de vitesse :
\begin{equation}
 c_{\rm int}=c_{\rm vac}=\widehat c\,u_C,\qquad
 \ell_0=c_{\rm int}t_0,\qquad
 c_{\rm int}=(\mu_{\rm vac}\varepsilon_{\rm vac})^{-1/2}.
 \label{cr28:eq:same-velocity}
\end{equation}
La dernière égalité est une identité de sections dans des bases dimensionnelles compatibles. La démonstration radiale du \cref{g26:sec:rigidez-radial-es} construit d'abord \(L=(5759/23040)\alpha^{16}L_*\), puis détermine
\begin{equation}
 G_{\rm ret}=\frac{c_{\rm int}^3L^2}{\hbar_{\rm ret}}
 =\frac{L^2}
 {\hbar_{\rm ret}(\mu_{\rm vac}\varepsilon_{\rm vac})^{3/2}}.
 \label{cr28:eq:constitutive-gravity}
\end{equation}
La seconde expression résulte de la substitution de l'identité constitutive dans la première ; elle n'introduit aucun nouvel étalonnage des deux canaux. En coordonnées, si \(c_{\rm int}=c_*\widehat c\), alors \(G_{\rm ret}=c_*^3\widehat c^{\,3}L^2/\hbar_{\rm ret}\). Le facteur de propagation n'intervient qu'une seule fois. La récupération de Catalan, le transport composé et la réponse radiale conservent ainsi la même origine angulaire, ses marques et sa normalisation.
% END CR28_CONSTITUTIVE_RADIAL_LINK

